\chapter{Technologies and Theoretical Elements Used}
\label{chap:technologies}

This chapter describes the specific technologies that ended up in
\nidsname{}'s stack and the reasoning behind each choice. The list is
long because the project spans backend, frontend, machine learning,
packet capture and packaging. Each section is self-contained and may
be read in any order.

\section{Python}
\label{sec:tech-python}

The whole backend runs on \textbf{Python}~\cite{python} 3.12. I
picked the version deliberately: the launcher refuses anything
newer than 3.13 because scikit-learn and numpy did not yet ship
Windows wheels for 3.14 at the time I started this work, which
would have forced a source build needing the Visual~C++ build
tools --- a seven-gigabyte detour every evaluator of this thesis
would have hit on a fresh Windows laptop. \file{run.py} checks the
interpreter version on the first banner line and prints an
upgrade path when the check fails; I added that gate the day I
watched a friend spend forty minutes trying to compile numpy from
source on 3.14.

Three reasons drove the language choice. The first is the
scientific Python stack. \code{from sklearn.ensemble import
RandomForestClassifier} on line one of a file is hard to beat
when the alternative is linking against a C++ library and
maintaining a binding layer. The second is asyncio. Python's
async primitives in 3.11+ are mature enough to host a real-time
event loop without dropping into threads everywhere, and FastAPI
is built directly on that substrate, so the entire packet
pipeline can stay in one coroutine context. The third is
pedagogical: Python is the language students at my faculty read
fluently, and a defence committee evaluating the thesis should
not have to learn a new syntax to follow the code.

The codebase uses type hints throughout (\code{from \_\_future\_\_ import
annotations}) and validated by static review rather than a strict type
checker. The annotations are not enforced at runtime by mypy because
the project does not currently require that discipline, but the
information is present for any future developer who wants to add it.

\section{FastAPI}
\label{sec:tech-fastapi}

The REST and WebSocket surface of \nidsname{} is served by
\textbf{FastAPI}~\cite{ramirez2022fastapi}, the asynchronous Python
framework Sebastián Ramírez published in 2018. Under the hood
FastAPI stacks on top of \textbf{Starlette} (the ASGI toolkit
that handles raw HTTP and WebSocket plumbing) and
\textbf{Pydantic} (the runtime validation library used for every
request and response body in the project).

I evaluated Flask, Django and aiohttp before settling on FastAPI.
What pushed me to commit was four properties, in roughly this
order. The dependency-injection system came first. Every
protected route in the codebase carries
\code{Depends(get\_current\_user)} on its signature, and the
framework handles JWT decoding, the database lookup and the
401 emission without a line of boilerplate inside the route
body --- one of the few cases in Python where a decorator
genuinely earns its keep. Then the auto-generated OpenAPI
documentation at \file{/docs}: zero configuration, every endpoint
browsable and testable from the Swagger UI in the browser. I used
this constantly during development to fire test requests at
the alerts filter and the WebSocket auth without ever writing
curl commands. Third, the Pydantic~v2 integration; request and
response bodies are declared as typed classes and the framework
moves JSON across the wire as Python objects with validation on
both sides. Fourth, the asynchronous-by-default story --- every
endpoint in \nidsname{} is an \code{async def} function and the
WebSocket handlers share the same event loop as the dispatcher,
which mattered more than I expected once the live packet stream
started filling up the loop.

FastAPI runs under \textbf{Uvicorn} in development, with the
\code{--reload} flag so saving a Python file restarts the server. In
production a process manager (gunicorn with the
\code{uvicorn.workers.UvicornWorker} class) would be more appropriate,
but the thesis scope does not include production deployment.

\section{SQLAlchemy 2 (Asynchronous)}
\label{sec:tech-sqlalchemy}

Persistence in \nidsname{} goes through
\textbf{SQLAlchemy~2}~\cite{sqlalchemy}, the major release that
unified the older session and core APIs around the
async-capable \code{AsyncSession}. The async story matters here
because the dispatcher and every API route are coroutines, and
mixing synchronous SQLAlchemy calls into that pipeline would
either block the event loop or force a thread pool I do not
want to maintain.

The data model is declared with the modern typed-mapping syntax in
\file{backend/app/models/}, one file per table. Each model uses
\code{Mapped[T]} annotations for SQLAlchemy 2 type inference. For
example, the \code{Alert} table is declared with fields like:

\begin{lstlisting}[style=pyfile,caption={Excerpt from backend/app/models/alert.py},label={lst:tech-alert-model}]
class Alert(Base):
    __tablename__ = "alerts"

    id:        Mapped[int]       = mapped_column(primary_key=True)
    ts:        Mapped[datetime]  = mapped_column(default=lambda: datetime.now(timezone.utc))
    severity:  Mapped[str]       = mapped_column(String(16))
    category:  Mapped[str]       = mapped_column(String(64))
    detection_method: Mapped[str] = mapped_column(String(16))
    ml_score:  Mapped[float|None] = mapped_column(Float, nullable=True)
    # ... and so on
\end{lstlisting}

The asynchronous engine is created once at boot in
\file{backend/app/core/database.py} and wrapped in an
\code{async\_sessionmaker} so every request can pull a fresh session.
The \code{get\_db} dependency yields the session and closes it in a
\code{finally} block, which guarantees no session leaks even if an
endpoint raises.

\section{aiosqlite and SQLite}
\label{sec:tech-sqlite}

Storage lives in \textbf{SQLite}~\cite{sqlite} talked to through
the asynchronous \textbf{aiosqlite} driver. Picking SQLite for a
SIEM-style application is the unusual call here --- the textbook
answer is PostgreSQL --- and I want to be honest about why I
made it. One overriding reason: deployment cost. SQLite is a
single file on disk. There is no separate server to install, no
\code{pg\_hba.conf} to configure, no backup cron, no port to
expose. For the actual deployment target of this thesis (a
laptop running an academic demo, plus the future small-network
template the codebase implies) the trade-off is correct. For a
production SOC swallowing thousands of alerts a minute SQLite
would be the wrong answer, and the SQLAlchemy abstraction makes
the swap to Postgres a one-line URL change.

SQLite has one architectural quirk worth mentioning: it serialises
writes (only one writer can hold the database lock at a time). Under
the synthetic traffic load used during development (24--40 packets
per second, with occasional attack bursts of 60--90 packets) this is
not a bottleneck. Even under the heaviest synthetic scenario (the
ICMP flood from the Attack Console, which generates upwards of
ninety packets per burst), the writes are batched at the alert
level rather than the packet level, so the load on SQLite is
still measured in tens of writes per second.

\section{Pydantic v2}
\label{sec:tech-pydantic}

The validation boundary between JSON on the wire and Python
objects in the route handlers is owned by
\textbf{Pydantic~v2}~\cite{pydantic}. Each request and response
schema lives in \file{backend/app/schemas/}, declared as a typed
class --- one model per resource, named after it
(\code{AlertOut}, \code{RuleIn}, \code{MLTrainRequest}). The
framework checks every payload at runtime and either coerces it
into the model or replies with a 422 carrying a precise error
trail down to the offending field, which removed about half the
error-handling boilerplate I would otherwise have written by
hand.

A small but important Pydantic detail: response models with
\code{model\_config = ConfigDict(from\_attributes=True)} can be built
from SQLAlchemy ORM instances directly. The endpoint code reads:

\begin{lstlisting}[style=pyfile]
@router.get("/alerts/{alert_id}", response_model=AlertOut)
async def get_alert(alert_id: int, db: AsyncSession = Depends(get_db)) -> AlertOut:
    row = (await db.execute(select(Alert).where(Alert.id == alert_id))).scalar_one_or_none()
    if not row:
        raise HTTPException(404, "Alert not found")
    return AlertOut.model_validate(row)
\end{lstlisting}

\noindent which is essentially as concise as the equivalent operation
in any modern Python ORM-based framework.

\section{scikit-learn}
\label{sec:tech-sklearn}

The machine-learning side of the project sits on top of
\textbf{scikit-learn}~\cite{pedregosa2011scikit}. I touch a small
slice of the library: \code{RandomForestClassifier} and
\code{IsolationForest} from \code{sklearn.ensemble}, plus four
metrics (\code{accuracy\_score}, \code{precision\_score},
\code{recall\_score}, \code{f1\_score}) from
\code{sklearn.metrics}. The training data is synthesised inside
the detector itself in
\file{backend/app/engine/ml\_detector.py} with numpy and a
seeded random generator
(\code{numpy.random.default\_rng(42)}) so that the same
training session reproduces the same model bit-for-bit on any
machine that runs the project. Trained artefacts are pickled
with \code{joblib.dump} into
\file{backend/data/models/*.joblib} and loaded lazily by
\code{joblib.load} when the operator activates them from the UI.

The two algorithms share a common \code{score(pkt) ->
(label, score)} interface, so the dispatcher does not need to know
which one is running. The \code{is\_anomaly} flag on the detector
state distinguishes the two cases at inference time: for Random
Forest the score is \code{predict\_proba()[0][1]}; for Isolation
Forest it is the negative of \code{decision\_function()[0]}.

\section{Scapy}
\label{sec:tech-scapy}

For live packet capture I use \textbf{Scapy}~\cite{scapy}. The
library wraps the platform's raw-socket facility (libpcap on
Linux and macOS, Npcap on Windows) and presents each captured
frame as a stack of Python objects, one per protocol layer. The
project hits this surface in two places.

First, the \code{LiveSniffer} class in
\file{backend/app/engine/sniffer.py} calls \code{sniff()} on a
background thread and converts each captured frame to the internal
\code{Packet} dataclass. The conversion extracts source and
destination IPs from the IP layer, source and destination ports from
the TCP or UDP layer (or zero for ICMP), TCP flags as a
comma-separated string (\code{"S,A"} rather than the raw integer
\code{0x12}), and the payload bytes from the \code{Raw} layer if
present.

Second, the project ships an optional BPF filter. The
\code{SNIFFER\_BPF\_FILTER} environment variable is applied at
capture time by \code{sniff()} so uninteresting traffic is
dropped before reaching the dispatcher. The expression follows
the standard tcpdump syntax (e.g.\ \code{"tcp or icmp"} or
\code{"port not 5353"}) and can be left empty to capture
everything.

Scapy's dependency on a platform packet driver is the source of most
deployment friction. On Linux the kernel exposes raw sockets to
processes that have the \code{CAP\_NET\_RAW} capability, so a
non-root user can capture if the binary is granted the capability
with \code{setcap}. On Windows the user must install Npcap (a
five-megabyte installer from \href{https://npcap.com}{npcap.com}) and
run the launcher with administrator privileges. The project handles
the missing-Scapy case gracefully: if the import fails, the
\code{TrafficSimulator} is used unconditionally regardless of the
\code{SIMULATE\_TRAFFIC} setting.

\section{NumPy}
\label{sec:tech-numpy}

\textbf{NumPy}~\cite{numpy} is the numerical library underlying
scikit-learn. \nidsname{} uses NumPy directly only in the feature
extraction path inside the ML detector, where each captured
\code{Packet} is converted to a six-dimensional \code{np.ndarray}
before being passed to \code{model.predict()}. The arrays are
constructed with \code{dtype=float} so that the dispatcher does not
incur a copy when scikit-learn validates the input shape.

\section{React 18}
\label{sec:tech-react}

The analyst console is a \textbf{React~18}~\cite{react}
single-page application, navigated by React Router. State lives
in two places: React Context (authentication and the two live
WebSocket streams, both mounted at the Layout level) and
component-local \code{useState} hooks for everything else. No
Redux, no Zustand, no MobX --- the data shape is small enough
that adding a state-management library would have been weight
for its own sake.

I picked React over Vue, Svelte and Angular for an unglamorous
reason: familiarity. The data flow here is the canonical SPA
loop (server pushes JSON, client renders it), and any of the
four frameworks would have produced an equivalent dashboard.
What tipped the decision was the depth of the React ecosystem
around the two pieces I knew I needed --- a charting library and
an icon set. Both happen to be Recharts and lucide-react below,
and both were one \code{npm install} away.

One React~18 feature the project actually leans on is concurrent
rendering through Suspense boundaries on the main routes. With
that in place the console stays responsive to clicks and
keystrokes even while a slow \code{fetch} is in flight against
the backend, which became visible the first time I hit
\code{/api/alerts?limit=500} on a hot SQLite database.

\section{Vite}
\label{sec:tech-vite}

The frontend's dev server and production bundler is
\textbf{Vite}~\cite{vite}, written by Evan You (the same author
behind Vue). For new React projects Vite has been steadily
displacing the older Create-React-App + webpack pair, and the
reason is its dev-time architecture: instead of bundling the
whole app on every save, Vite serves each module as a separate
HTTP request and lets the browser's own ES-module loader chase
the dependency graph. The practical consequence is hot reloads
under a hundred milliseconds even on a project the size of the
analyst console. I noticed the difference the first time I came
back from working on a webpack project --- it stops feeling like
``save and wait'' and starts feeling like ``save and look''.

The Vite configuration for this project is minimal: a single
\code{vite.config.js} file with the React plugin and a proxy entry
that forwards \file{/api/} and \file{/ws/} to the FastAPI backend on
port 8000 during development. Production builds emit a static
\file{dist/} directory which could be served by any static-file web
server.

\section{TailwindCSS}
\label{sec:tech-tailwind}

The visual identity of the console is delivered by
\textbf{TailwindCSS}~\cite{tailwind}. The design language is
narrow on purpose: a deep navy background (\code{0A1220}), a
cyan primary accent (\code{22D3EE}), and three semantic
highlights --- fuchsia for ML-related state, amber for warnings
and attack-console interactions, and red for critical alerts.
Tailwind's utility-class approach lets me compose that palette
directly in the JSX markup; there is no parallel
\code{*.module.css} hierarchy to keep in sync with the
component tree.

I extended the default palette in \file{tailwind.config.js} with
three custom navy shades (\code{navy-850}, \code{navy-900},
\code{navy-950}) so the components can refer to named tokens
instead of raw hex strings. The final production stylesheet
weighs about twenty kilobytes gzipped --- small enough that the
console feels instantaneous over any network worth measuring.

\section{Recharts}
\label{sec:tech-recharts}

\textbf{Recharts}~\cite{recharts} is the React-native charting
library that draws the four charts on the dashboard: the live
packet-rate area chart, the severity distribution pie, the alert
timeline bar chart and the detection-method pie. The library renders
SVG under the hood and supports declarative composition --- a chart
is a tree of React components rather than a single black-box
configuration object, which makes it natural to override colours and
tooltip styling per-chart.

A small subtlety I had to resolve was the default chart animation,
which caused the live packet-rate area to wiggle every second as
Recharts interpolated between the old and new data series. Setting
\code{isAnimationActive=\{false\}} on the \code{Area} component, plus
\code{type="linear"} and \code{dot=\{false\}}, produced the stable
left-shifting curve I wanted.

\section{lucide-react}
\label{sec:tech-lucide}

\textbf{lucide-react} is the icon library --- a React port of the
Lucide icon set, itself a community fork of the original Feather
Icons. The icons are SVG-based, tree-shakeable (only the icons
actually imported are included in the final bundle) and visually
consistent because they share a single stroke style. Every button,
chip and badge in the console uses lucide icons.

\section{date-fns}
\label{sec:tech-datefns}

\textbf{date-fns}~\cite{datefns} formats the timestamps shown in the
UI. It was preferred over Moment.js (now in maintenance mode) and
Day.js (a Moment-compatible alternative) because date-fns is purely
functional, tree-shakeable, and the single \code{format} import is
enough for every place a timestamp appears.

\section{WebSocket Protocol}
\label{sec:tech-websocket}

\textbf{WebSocket}~\cite{ietf6455} (RFC 6455) is the bidirectional
streaming protocol that delivers live updates from the backend to
the analyst console. The two endpoints are \file{/ws/alerts} (each
new alert is pushed as a single JSON message) and \file{/ws/traffic}
(packet metadata batched in groups of up to 50 every 500\,ms to
keep the dashboard chart fluid without overwhelming the
network).

Authentication is handled through a token in the query string
because browser WebSocket clients cannot set custom request
headers --- this is a deliberate restriction of the standard. The
token is the same JWT used by the REST endpoints, and the
\code{ws/alerts.py} handler decodes it before subscribing the
connection to the alert bus.

\section{JWT and bcrypt}
\label{sec:tech-auth}

The authentication subsystem is built on two pieces.
\textbf{JSON Web Tokens (JWT)}~\cite{jwt} encode the signed-in user's
identity and role; the backend signs them with the HS256 symmetric
algorithm using the secret stored in
\file{backend/.env}. \textbf{bcrypt}~\cite{bcrypt} hashes user
passwords before they are stored. I deliberately removed the older
\code{passlib} library from the dependencies because passlib's
maintenance has lapsed and it does not work cleanly with the modern
bcrypt 4.x and 5.x releases that probe the \code{\_\_about\_\_.\_\_version\_\_}
attribute. Calling \code{bcrypt.hashpw} and \code{bcrypt.checkpw}
directly from \file{backend/app/core/security.py} avoids that whole
problem.

\section{Joblib}
\label{sec:tech-joblib}

\textbf{joblib}~\cite{joblib} is the persistence layer for the
trained machine-learning models. It is a thin wrapper around pickle
with optimised handling for large NumPy arrays --- the kind of
data structure a Random Forest with 200 trees produces. Each
trained model is saved to \file{backend/data/models/} with a name
that includes the timestamp, so re-training never overwrites a
previous artefact.

\section{PyYAML}
\label{sec:tech-yaml}

\textbf{PyYAML}~\cite{pyyaml} parses the default rule file. The
seed rules live in
\file{backend/data/rules/default\_rules.yaml} and are loaded into the
database the first time the application starts. The format is
straightforward enough that an operator can edit the file by hand
without learning the JSON match specification --- the YAML schema and
the database match column are identical apart from the surface
syntax.

\section{Uvicorn}
\label{sec:tech-uvicorn}

\textbf{Uvicorn}~\cite{uvicorn} is the ASGI server that runs the
FastAPI application. It is the reference implementation of the
ASGI specification, written by Tom Christie, and has become the
default choice for asynchronous Python web services. The launcher
invokes it as \code{python -m uvicorn app.main:app --host 0.0.0.0 --port 8000 --reload};
the \code{--reload} flag watches the file system and restarts the
server when source files change, which is invaluable during
development.

\section{Node.js and npm}
\label{sec:tech-node}

\textbf{Node.js} is the JavaScript runtime that powers the Vite
development server and the build pipeline. The project pins no
specific Node version but requires 18 LTS or newer, which is what
the major distributions and the official Node installer ship today.
\textbf{npm} is the package manager that installs the frontend
dependencies; it is bundled with the Node installer, so a user who
follows the standard installation procedure has both tools
available with no extra configuration.

\section{Operating System Considerations}
\label{sec:tech-os}

The project targets three operating systems: Windows 10/11, macOS
12+, and any reasonably recent Linux distribution. The
cross-platform launcher (\file{run.py}) detects the host operating
system via \code{platform.system()} and adjusts the venv path
(\file{Scripts/python.exe} on Windows versus \file{bin/python} on
Unix), the signal handling
(\code{CTRL\_BREAK\_EVENT} on Windows versus \code{SIGTERM} on Unix)
and the dependency installation procedure. The two wrappers
(\file{run.bat} for Windows, \file{run.sh} for Unix) handle the
shell-specific quirks --- mostly around how to invoke Python with
the right version --- and then delegate to \file{run.py} for the
actual logic.
